% Source: https://artofproblemsolving.com/community/c3h2126505_solve_this_as_fast_as_u_can
% Stable identity: aops-topic-2126505-q01
Call the quadrilateral $ABCD$ and the centre of the circle $I$.

\textbf{The other two sides also add up to $12$.} The two tangent segments from a vertex to the circle have equal length; call these lengths $t_A,t_B,t_C,t_D$. Each side consists of two tangent segments, for example $AB=t_A+t_B$. Hence
\[AB+CD=t_A+t_B+t_C+t_D=BC+DA,\]
so the perimeter is $2\cdot12=24$ cm. (Which pair of opposite sides is meant does not matter.)

\textbf{Area.} Join $I$ to the four vertices. This splits $ABCD$ into the triangles $IAB$, $IBC$, $ICD$ and $IDA$. The radius to a point of tangency is perpendicular to the side, so each of these triangles has height $2$ on its side of the quadrilateral.

\begin{asy}
usepackage("amsmath");
size(5.5cm);
defaultpen(fontsize(10pt));
real r=2; real[] dl={60,149.1307,40.8693,110}; real[] th; th[0]=-90;
for(int k=1;k<4;++k) th[k]=th[k-1]+dl[k-1];
pair[] T; for(int k=0;k<4;++k) T[k]=r*dir(th[k]);
pair[] V; for(int k=0;k<4;++k){ real m=th[k]+dl[k]/2; V[k]=(r/Cos(dl[k]/2))*dir(m); }
// V[k] lies between T[k] and T[k+1]
pair I=(0,0);
draw(V[0]--V[1]--V[2]--V[3]--cycle,linewidth(1));
draw(circle(I,r),blue);
for(int k=0;k<4;++k){ draw(I--V[k],gray+dashed); draw(I--T[k],red); dot(T[k],red); }
dot(I); label("$I$",I,dir(-120),fontsize(9pt));

label("$C$",V[2],dir(th[2]+dl[2]/2)); label("$D$",V[3],dir(th[3]+dl[3]/2));
label("$2$",T[0]/2,W,red+fontsize(8pt));
label("$A$",V[0],dir(th[0]+dl[0]/2)); label("$B$",V[1],dir(th[1]+dl[1]/2));
label("\mbox{}",(min(currentpicture,user=true)+max(currentpicture,user=true))/2,black); // keeps the text after the figure black when LaTeX typesets the labels inline
\end{asy}

Therefore
\[[ABCD]=\dfrac{1}{2}\cdot2\cdot(AB+BC+CD+DA)=\dfrac{1}{2}\cdot2\cdot24=\boxed{24}\text{ cm}^2 .\]
(In general the area of a polygon with an inscribed circle is the inradius times the semiperimeter.)
